% CVPR 2026 Paper Template; see https://github.com/cvpr-org/author-kit

\documentclass[10pt,twocolumn,letterpaper]{article}

%%%%%%%%% PAPER TYPE  - PLEASE UPDATE FOR FINAL VERSION
% \usepackage{cvpr}              % To produce the CAMERA-READY version
\usepackage[review]{cvpr}      % To produce the REVIEW version
% \usepackage[pagenumbers]{cvpr} % To force page numbers, e.g. for an arXiv version

% Import additional packages in the preamble file, before hyperref
\input{preamble}
% Local build shim for Tectonic/XeTeX.
%
% The official CVPR style requests PSNFSS Times/Helvetica fonts (`ptm`/`phv`).
% Tectonic's XeTeX-based LaTeX defaults to TU encoding here, where those legacy
% PSNFSS families are not defined. Force the classic T1 font encoding so the
% official Times-like PSNFSS fonts and their bold shapes are used instead of
% silently falling back to non-bold Computer Modern or unrelated system fonts.
\usepackage{iftex}
\ifXeTeX
  \usepackage[T1]{fontenc}
\fi


\definecolor{cvprblue}{rgb}{0.21,0.49,0.74}
\usepackage[pagebackref,breaklinks,colorlinks,allcolors=cvprblue]{hyperref}

%%%%%%%%% PAPER ID  - PLEASE UPDATE
\def\paperID{*****}
\def\confName{CVPR}
\def\confYear{2026}

\newcommand{\benchmark}{CEPO-Probe}
\newcommand{\dataset}{ClaimEvidence-6K}
\newcommand{\answerdpo}{CEPO Answer-DPO}
\newcommand{\cepodual}{CEPO-Dual}

%%%%%%%%% TITLE - PLEASE UPDATE
\title{\benchmark{}: A Claim-Evidence Consistency Benchmark for Diagnosing Vision-Language Preference Tuning}

%%%%%%%%% AUTHORS - PLEASE UPDATE
\author{Anonymous \confName{} Submission}

\begin{document}
\maketitle

\begin{abstract}
Preference tuning can improve the short answers of vision-language models
(VLMs), but a correct yes/no answer does not imply that the model has verified
the right visual evidence. We introduce \benchmark{}, a compact diagnostic
benchmark for claim-evidence consistency. The benchmark is built from
\dataset{}, a 6k-row collection of object, attribute, relation, and
wrong-evidence preference examples derived from COCO and GQA annotations with
zero train/eval image overlap against the fixed evaluation suite. Unlike
ordinary answer-level evaluation, \benchmark{} asks a model to judge whether a
candidate visual claim is supported by a specific evidence label or relation,
including same-answer cases where only the evidence is wrong. We evaluate
Qwen2.5-VL-7B-Instruct, an answer-only DPO baseline, and \cepodual{}, which
mixes answer preferences with direct evidence-verifier preferences. The
answer-only baseline improves short-answer accuracy but increases yes bias on
GQA, while \cepodual{} preserves the short-answer gains and improves
wrong-evidence rejection from 34.3\% to 70.2\%. The improvement is uneven:
object and attribute mismatches become much easier to reject, but relation
reversals remain fragile. External POPE and AMBER checks stay close to the
answer-only baseline, so the supported claim is diagnostic rather than a broad
hallucination-benchmark win. These results show that short-answer preference
tuning and evidence consistency are separable and should be reported together.
\end{abstract}

\section{Introduction}

Vision-language models can answer simple visual questions fluently while using
the wrong visual support. For example, a model may answer that a sink is
present but cite a potted plant as the evidence, or accept that a glass is left
of a grape while the provided relation swaps the two objects. In such cases,
answer correctness and evidence correctness are related but not identical. The
distinction is especially important for preference tuning, where a chosen
answer is often preferred without checking whether it is tied to the correct
visual evidence.

Direct Preference Optimization (DPO)~\cite{dpo} is attractive for VLM
alignment because it trains from paired preferences without a separate reward
model. However, an answer-only preference pair can be satisfied by learning a
better response prior or a stronger yes/no bias. Even evidence-flavored
training responses may fail to change the behavior we care about if evaluation
only asks for a short answer. We therefore ask a diagnostic question: when
preference tuning improves ordinary answers, does it also improve the model's
ability to verify claim-evidence consistency?

We introduce \benchmark{}, a small benchmark designed to isolate this question.
Each example contains an image, a question, a candidate claim, and a candidate
evidence item. The model must return a structured judgment indicating whether
the evidence supports the claim. The key split is wrong-evidence verification:
examples where the answer text can remain plausible, but the candidate
evidence is mismatched. This forces evaluation to look beyond the final answer
and ask whether the model accepts the correct object, attribute, or relation.

Our goal is not to claim that a new preference-tuning method solves visual
hallucination. Instead, we use \benchmark{} as a diagnostic instrument. We
compare a base model, \answerdpo{}, and \cepodual{} under the same backbone,
data source, and evaluation suite. The results support a useful but bounded
conclusion: \cepodual{} improves explicit evidence verification while leaving
relation-level mismatches difficult, so evidence consistency should be measured
directly rather than inferred from answer accuracy. We therefore do not frame
\cepodual{} as a general external-benchmark winner; its main value is exposing
and improving a failure mode that ordinary short-answer metrics can miss.

Our contributions are:
\begin{itemize}
    \item We define \benchmark{}, a compact claim-evidence consistency
    benchmark that separates answer correctness from evidence correctness.
    \item We construct \dataset{}, a 6k-row controlled set with object,
    attribute, relation, and wrong-evidence examples plus fixed supported and
    wrong-evidence probe splits.
    \item We provide a deterministic evaluation protocol for VLM preference
    tuning that reports short-answer transfer, support accuracy,
    evidence-label extraction, wrong-evidence rejection, slice breakdowns, and
    invalid-format rates.
    \item We show that \answerdpo{} improves short-answer behavior without
    reliably improving wrong-evidence rejection, while \cepodual{} improves the
    explicit probe but still struggles on relation reversals and remains close
    to \answerdpo{} on POPE/AMBER sanity checks.
\end{itemize}

\section{Related Work}

\paragraph{Object hallucination and VLM evaluation.}
Object hallucination has been studied by checking generated captions or
answers against visual annotations. CHAIR~\cite{chair} measures hallucinated
objects in captions, while POPE~\cite{pope} recasts object hallucination as
polling-style yes/no questions. \benchmark{} follows this controlled
evaluation spirit, but asks whether the model connects a claim to the evidence
that supports it.

\paragraph{Preference tuning.}
DPO and related preference-optimization methods align models from
chosen/rejected examples. In VLM settings, answer-level preferences can improve
final responses while leaving the grounding signal implicit. Our benchmark
tests whether such tuning changes evidence verification, not only answer
selection.

\paragraph{Grounding and rationales.}
Grounded VLM training can use boxes, masks, region descriptions, or
human-written rationales. These signals are valuable but costly. \benchmark{}
uses weaker annotation-derived evidence labels and relations, allowing a small
diagnostic study that can be regenerated from COCO~\cite{coco} and
GQA~\cite{gqa}-style resources.

\section{\benchmark{} Benchmark}

\subsection{Task Definition}

Each \benchmark{} example contains an image $I$, a visual question $q$, a
candidate claim $c$, and candidate evidence $e$. The model receives the image
and a prompt such as:
\begin{quote}
\small
\textit{Visual question: Is there a sink in the image? Candidate claim: sink
visible. Candidate evidence: label=potted plant. Decide whether the candidate
evidence supports the claim in the image.}
\end{quote}
It should return JSON with \texttt{answer}, \texttt{claim},
\texttt{evidence\_label}, and \texttt{support}. The support label is
\texttt{supported}, \texttt{contradicted}, or \texttt{wrong\_evidence}. The
headline metric is wrong-evidence rejection, where both
\texttt{wrong\_evidence} and \texttt{contradicted} count as rejecting a
mismatched evidence item.

\subsection{Data Construction}

\dataset{} is derived from COCO and GQA annotations. COCO contributes
object-existence claims with visible object labels and absent-object negatives.
GQA contributes color/material attributes and left/right spatial relations from
scene graphs. We also create same-answer wrong-evidence rows where the answer
text is not the discriminating signal: the candidate evidence is a wrong object
label, mismatched attribute, or reversed relation.

\begin{table}[t]
\centering
\small
\begin{tabular}{lrr}
\toprule
Canonical slice & Rows & Negative type \\
\midrule
Object existence & 2{,}000 & absent/wrong object \\
Attribute color/material & 1{,}500 & attribute mismatch \\
Left/right relation & 1{,}500 & relation reversal \\
Same-answer wrong evidence & 1{,}000 & evidence mismatch \\
\midrule
Total & 6{,}000 & -- \\
\bottomrule
\end{tabular}
\caption{\dataset{} composition. The source mix is 2{,}390 COCO-derived rows
and 3{,}610 GQA-derived rows.}
\label{tab:data}
\end{table}

The supported-evidence probe contains 600 rows: 162 object, 276 attribute, and
162 relation examples. The wrong-evidence probe contains 400 rows: 126 object,
149 attribute, and 125 relation examples. The canonical data check reports no
blocking errors, no identical \answerdpo{} chosen/rejected rows, 1{,}000
same-answer wrong-evidence rows, and zero train/eval image overlap against
COCO held-out, GQA simple, Hard COCO, Base-error-mined, POPE, and AMBER evals.
A 200-row annotation-consistency audit of \dataset{} passes all sampled checks:
200/200 chosen answers match their supported claims, 170/170 applicable
rejected answers match contradicted claims, 200/200 chosen evidence entries
match the chosen claims, and 200/200 rejected errors match the intended
negative type. Same-answer wrong-evidence rows mark rejected-answer wrongness
as not applicable because the answer is intentionally plausible.

\paragraph{Data card.}
The benchmark is intended for diagnostic evaluation of VLM grounding and
preference tuning. It is derived from COCO and GQA annotations, while POPE and
AMBER are used only as external evaluation sources. Image redistribution is
not required by the benchmark artifacts in this repository; examples store
local image paths and annotation-derived labels. Known limitations include
annotation incompleteness, weak textual evidence labels, missing boxes for some
GQA-derived rows, and narrow coverage of simple object, attribute, and
left/right relation claims.

\section{Evaluation Protocol}

\subsection{Model Groups}

The main comparison uses three groups:
\begin{itemize}
    \item \textbf{Base Instruct:} Qwen2.5-VL-7B-Instruct~\cite{qwen25vl}
    without preference tuning.
    \item \textbf{\answerdpo{}:} a LoRA-DPO baseline trained for one epoch on
    6{,}000 short-answer preference rows.
    \item \textbf{\cepodual{}:} a LoRA-DPO baseline trained for one epoch on
    the same 6{,}000 answer rows plus 2{,}000 evidence-verifier preference
    rows, split evenly between supported and wrong-evidence targets.
\end{itemize}
Both trained models use the same backbone, DPO $\beta=0.1$, LoRA rank 16, and
ZeRO-2 training setup. We avoid additional ablations in the main paper to keep
the benchmark diagnosis focused.

\subsection{Metrics}

Short-answer transfer uses COCO object existence, GQA simple, Hard COCO, and a
Base-error-mined diagnostic set. Metrics are accuracy, F1, false-positive rate
(FPR), false-negative rate (FNR), yes rate, and invalid/other rate.

\benchmark{} uses supported-evidence and wrong-evidence probes. We report
support accuracy, wrong-evidence rejection accuracy, evidence-label accuracy,
claim-match rate, slice-level wrong-evidence rejection, and invalid JSON rate.
Evidence-label accuracy is treated as a format/extraction check, since a model
can copy the candidate label without verifying that it is correct. The primary
diagnostic is therefore the support decision, especially on wrong-evidence
rows.

\section{Experiments}

\subsection{Short-Answer Transfer}

\begin{table*}[t]
\centering
\footnotesize
\begin{tabular}{@{}lccccccc@{}}
\toprule
Model & COCO & COCO FPR/FNR & GQA & GQA FPR/FNR & Hard & Hard FPR/FNR & BEM \\
\midrule
Base Instruct & 96.0 & \textbf{1.6} / 6.4 & 76.2 & \textbf{14.4} / 33.2 & 94.4 & \textbf{3.8} / 7.4 & 0.0 \\
\answerdpo{} & \textbf{97.0} & 1.8 / \textbf{4.2} & 76.8 & 23.6 / 22.8 & \textbf{95.0} & 5.0 / \textbf{5.0} & 18.0 \\
\cepodual{} & \textbf{97.0} & 1.8 / \textbf{4.2} & \textbf{77.1} & 24.4 / \textbf{21.4} & 94.9 & 5.2 / \textbf{5.0} & \textbf{19.4} \\
\bottomrule
\end{tabular}
\caption{Short-answer transfer under the normal yes/no prompt. Values are
percentages. \answerdpo{} and \cepodual{} improve recall and accuracy, but
both increase yes bias on GQA and Hard COCO relative to the base model.}
\label{tab:short}
\end{table*}

Table~\ref{tab:short} shows that answer-level DPO changes ordinary yes/no
behavior. \answerdpo{} improves COCO accuracy from 96.0\% to 97.0\% and Hard
COCO from 94.4\% to 95.0\%, mainly by lowering FNR. The trade-off is a larger
GQA false-positive rate: 23.6\% for \answerdpo{} versus 14.4\% for the base
model. \cepodual{} preserves the COCO gain, slightly improves GQA accuracy,
and nearly matches Hard COCO, but it does not remove the short-answer yes-bias
trade-off. This motivates the explicit evidence probe: answer accuracy alone
does not tell us whether the model verifies the proposed evidence.

\subsection{Claim-Evidence Consistency}

\begin{table*}[t]
\centering
\small
\begin{tabular}{lcccccc}
\toprule
Model & Supp. Acc & Wrong Rej. & Obj Wrong & Attr Wrong & Rel. Wrong & Invalid JSON \\
\midrule
Base Instruct & 86.5 & 44.3 & 59.5 & 53.7 & 17.6 & 0.0 \\
\answerdpo{} & 90.7 & 34.3 & 45.2 & 45.0 & 10.4 & 0.0 \\
\cepodual{} & \textbf{95.2} & \textbf{70.2} & \textbf{96.0} & \textbf{85.9} & \textbf{25.6} & 0.0 \\
\bottomrule
\end{tabular}
\caption{\benchmark{} probe results. Supp. Acc is support accuracy on the
600-row supported-evidence probe. Wrong Rej. is rejection accuracy on the
400-row wrong-evidence probe; the next three columns break that rejection rate
down by base task type. Values are percentages.}
\label{tab:probe}
\end{table*}

Table~\ref{tab:probe} exposes the separation between answer preference and
evidence verification. \answerdpo{} improves supported-evidence support
accuracy over the base model, but its wrong-evidence rejection drops from
44.3\% to 34.3\%. It often accepts plausible candidate evidence even when the
evidence label or relation is mismatched. \cepodual{} changes this behavior:
wrong-evidence rejection rises to 70.2\%, a 36.0 point gain over \answerdpo{},
while supported-evidence support accuracy also improves to 95.2\%.

The slice breakdown is more cautious. \cepodual{} strongly improves wrong
object evidence (96.0\%) and wrong attribute evidence (85.9\%), but relation
reversals remain hard at 25.6\%. This means the positive result should not be
read as solving evidence grounding. Rather, \benchmark{} identifies where
direct verifier preferences help and where a model still fails.

\subsection{External Sanity Checks}

On POPE random/popular/adversarial, \cepodual{} reaches 89.8\%, 88.7\%, and
87.3\% accuracy, respectively, compared with 89.6\%, 88.5\%, and 87.1\% for
\answerdpo{}. On AMBER discriminative, \cepodual{} reaches 88.2\% accuracy
versus 88.2\% for \answerdpo{} after rounding. We treat these as sanity
checks, not as the main source of improvement. They show that adding verifier
rows does not cause a large out-of-distribution short-answer collapse, but the
differences from \answerdpo{} are small. They also mirror the internal
short-answer trade-off: the trained models lower FNR while raising FPR
relative to Base Instruct.

\section{Discussion}

\benchmark{} is designed to reveal a failure mode hidden by ordinary answer
accuracy. \answerdpo{} improves short-answer behavior, but the explicit probe
shows that this does not imply stronger rejection of mismatched evidence. In
fact, \answerdpo{} has lower wrong-evidence rejection than the base model.
\cepodual{} demonstrates that direct evidence-verifier preferences can improve
the measured evidence behavior without sacrificing the short-answer gains, but
the relation slice remains weak. A useful next benchmark extension would add
region-aware relation checks instead of relying only on textual relation
labels.

The results also clarify how to evaluate evidence-aware preference tuning. If
a model is trained to reason about claims but tested only with short answers,
improvements can be invisible or confused with yes/no bias. If it improves
structured evidence probes while failing ordinary answers, it is better
understood as an evidence verifier than as a general answerer. Reporting both
views prevents overclaiming. The synchronized claim of this paper is therefore
not that \cepodual{} solves object hallucination or dominates POPE/AMBER; it is
that explicit verifier preferences can improve a separately measured
claim-evidence consistency skill that answer-only DPO can weaken.

\section{Limitations}

\benchmark{} is intentionally small and controlled. It is derived from COCO
and GQA annotations rather than newly collected human evidence labels. Textual
evidence labels are weaker than boxes, masks, or human rationales, and the
rule-based scorer cannot recognize every semantically valid free-form
description. Some GQA-derived rows lack boxes, so the current probe cannot
check region localization. The experiment uses one 7B backbone and one
small-data LoRA-DPO recipe; broader model families and seeds are needed before
drawing general conclusions about VLM alignment. Finally, the relation results
show that accepting or rejecting textual relation evidence is still far from a
robust spatial grounding test.

\section{Conclusion}

We presented \benchmark{}, a diagnostic benchmark for claim-evidence
consistency in VLM preference tuning. The benchmark separates ordinary
short-answer performance from the ability to verify whether a visual claim is
supported by the correct evidence. In a controlled three-group experiment,
\answerdpo{} improves short-answer accuracy while weakening wrong-evidence
rejection, whereas \cepodual{} preserves the answer gains and substantially
improves explicit evidence consistency. External POPE/AMBER results remain
close to the answer-only baseline, so the result should be read as a diagnostic
separation rather than a general hallucination solution. The remaining relation
failures are the point of the benchmark: short-answer gains alone can hide
persistent wrong-evidence errors, so evidence consistency should be measured
directly.

{
    \small
    \bibliographystyle{ieeenat_fullname}
    \bibliography{references}
}

% Optional supplementary material. Compile `main_full.tex` or run `make full`
% to include the appendix without changing the main review PDF.
\ifdefined\cvprincludeappendix
  \clearpage
  \appendix
  \section{Supplementary Material}
\captionsetup{hypcap=false}

This appendix collects details that support the main \benchmark{} claim without
expanding the primary review paper beyond the five-group comparison. The
sections are organized by purpose: reproducibility setup, additional
quantitative checks, and qualitative/scoring details. The appendix is not meant
to introduce a second story; it records the controls behind the same bounded
conclusion in the main paper.

\subsection{Reproducibility Setup and Artifacts}

All trained variants use the same backbone, optimizer objective, adapter
configuration, and distributed setup. The only experimental variable is the
composition of preference rows: answer rows, verifier rows, or their mixture.
This control is important because the paper interprets differences between
groups as evidence about supervision format rather than architecture,
optimizer, or decoding changes.

\begin{center}
\begin{minipage}{0.58\linewidth}
\centering
\footnotesize
\begin{tabular}{ll}
\toprule
Setting & Value \\
\midrule
Backbone & Qwen2.5-VL-7B-Instruct \\
Training objective & Direct Preference Optimization \\
Adapter & LoRA \\
Epochs & 1 \\
DPO beta & 0.1 \\
LoRA rank & 16 \\
Answer rows & 0 or 6{,}000 \\
Verifier rows & 0, 500, 1{,}000, or 2{,}000 \\
Distributed setup & ZeRO-2 \\
\bottomrule
\end{tabular}
\captionof{table}{Shared training setup for the preference-tuned groups. The
main paper reports five groups; the 500-verifier-row run is retained only as
an auxiliary sensitivity check.}
\label{tab:app_train}
\end{minipage}
\end{center}

The completed artifacts are small enough to inspect and regenerate. The
canonical \dataset{} file contains 6{,}000 rows. The selected \dualtwok{} DPO
export contains 8{,}000 rows: 6{,}000 answer-preference rows plus 2{,}000
verifier rows. The evidence probes are fixed at 600 supported-evidence rows
and 400 wrong-evidence rows. This fixed-probe design prevents the ablation from
changing both the training data and the evaluation set at the same time.

Lightweight local regeneration covers data checking, audit-sheet completion,
CEPO-Dual export, and LLaMA-Factory data preparation. Heavy training and full
inference are submitted through Slurm, following the repository rule that GPU
jobs should not be run directly in the interactive environment. The relevant
Slurm entry points are:
\begin{quote}
\footnotesize
\texttt{submit\_cepo\_dual\_pipeline.sh}\\
\texttt{submit\_cepo\_ablation\_pipeline.sh}
\end{quote}
The paper numbers are then refreshed with the deterministic scoring and
summary scripts:
\begin{quote}
\footnotesize
\texttt{score\_object\_eval.py}\\
\texttt{score\_cepo\_evidence\_probe.py}\\
\texttt{bootstrap\_cepo\_probe\_ci.py}\\
\texttt{summarize\_cepo\_ablation\_results.py}
\end{quote}

\subsection{Additional Quantitative Checks}

The main paper keeps the primary comparison to five groups, as required by the
submission plan. Table~\ref{tab:app_ablation} includes the auxiliary
CEPO-Dual-500 run that was useful during analysis but intentionally omitted
from the main body. The table supports two readings. First, verifier-only
supervision improves wrong-evidence rejection but does not replace answer
preference tuning for ordinary short-answer transfer. Second, the mixed
CEPO-Dual variants need enough verifier data before the wrong-evidence probe
clearly improves: CEPO-Dual-500 and CEPO-Dual-1k are weaker than the selected
CEPO-Dual-2k setting.

\begin{center}
\begin{minipage}{\linewidth}
\centering
\scriptsize
\setlength{\tabcolsep}{3.5pt}
\begin{tabular}{@{}lrrrrrrrrr@{}}
\toprule
Setting & Ans. rows & Ver. rows & COCO & GQA & Hard & Supp. & Wrong Rej. & Obj/Attr/Rel & Parse Fail \\
\midrule
Base Instruct & 0 & 0 & 96.0 & 76.2 & 94.4 & 86.5 & 44.3 & 59.5 / 53.7 / 17.6 & 0.0 \\
\answerdpo{} & 6{,}000 & 0 & 97.0 & 76.8 & 95.0 & 90.7 & 34.3 & 45.2 / 45.0 / 10.4 & 0.0 \\
\evidencedpo{} & 0 & 2{,}000 & 95.9 & 76.7 & 94.4 & 87.8 & 53.0 & 77.0 / 63.8 / 16.0 & 0.0 \\
CEPO-Dual-500 & 6{,}000 & 500 & 96.9 & 76.8 & 95.0 & 91.3 & 39.5 & 56.3 / 50.3 / 9.6 & 0.0 \\
\dualonek{} & 6{,}000 & 1{,}000 & 97.0 & 77.2 & 94.9 & 93.2 & 48.5 & 74.6 / 59.1 / 9.6 & 0.0 \\
\dualtwok{} & 6{,}000 & 2{,}000 & 97.0 & 77.1 & 94.9 & 95.2 & 70.3 & 96.0 / 85.9 / 25.6 & 0.0 \\
\bottomrule
\end{tabular}
\captionof{table}{Full verifier-count ablation from the released artifacts.
The main paper omits CEPO-Dual-500 to keep the primary comparison to at most
five groups.}
\label{tab:app_ablation}
\end{minipage}
\end{center}

The ablation also shows why the paper does not present CEPO-Dual as a broad
hallucination method. The largest gain is on the internal evidence probe, not
on every external yes/no benchmark. Table~\ref{tab:app_external} therefore
serves as a sanity check. CEPO-Dual-2k stays close to CEPO Answer-DPO on POPE
and AMBER, which means the verifier rows do not cause an obvious collapse in
ordinary polling behavior. The differences are small, so the external table is
supporting evidence for stability rather than the source of the main claim.

\begin{center}
\begin{minipage}{0.55\linewidth}
\centering
\footnotesize
\begin{tabular}{lccc}
\toprule
Eval & Base & \answerdpo{} & \dualtwok{} \\
\midrule
POPE random & 88.7 & 89.6 & \textbf{89.8} \\
POPE popular & 87.7 & 88.5 & \textbf{88.7} \\
POPE adversarial & 86.7 & 87.1 & \textbf{87.3} \\
AMBER discr. & 87.6 & \textbf{88.2} & \textbf{88.2} \\
\bottomrule
\end{tabular}
\captionof{table}{External yes/no sanity checks. Values are accuracy
percentages under the same deterministic parser.}
\label{tab:app_external}
\end{minipage}
\end{center}

\subsection{Qualitative Analysis and Scoring Rules}

The most persistent failure mode is relation verification. Object and
attribute evidence can often be checked by recognizing one label or property,
but a left/right relation requires the model to preserve subject, object, and
direction at the same time. Table~\ref{tab:app_qual} gives representative
probe outputs. The successful wrong-relation case shows that CEPO-Dual-2k can
learn to reject swapped evidence, while the failure case shows the remaining
weakness: if both objects are visible, the model may still accept the swapped
relation as supported.

\begin{center}
\begin{minipage}{\linewidth}
\centering
\scriptsize
\begin{tabular}{@{}p{0.14\linewidth}p{0.34\linewidth}p{0.16\linewidth}p{0.13\linewidth}p{0.10\linewidth}@{}}
\toprule
Type & Candidate claim and evidence & Target & \dualtwok{} output & Reading \\
\midrule
Supported attr. &
claim: vehicle is white; evidence: label=vehicle, attribute=white &
supported & supported & correct \\
Supported material &
claim: window is glass; evidence: label=window, attribute=glass &
supported & supported & correct \\
Wrong object &
claim: sink visible; evidence: label=potted plant &
wrong evidence & wrong evidence & correct \\
Wrong relation &
claim: glass left\_of grape; evidence swaps subject/object &
wrong evidence & contradicted & correct rejection \\
Wrong relation &
claim: man right\_of parking meter; evidence swaps subject/object &
wrong evidence & supported & relation failure \\
\bottomrule
\end{tabular}
\captionof{table}{Representative probe outputs. Examples are anonymized to
local image references and show why relation reversals remain the hardest
slice.}
\label{tab:app_qual}
\end{minipage}
\end{center}

Scoring follows the same conservative principle. Short-answer metrics use a
deterministic yes/no parser. Outputs with explicit uncertainty such as
``cannot determine'' or ``not sure'' are refusals; outputs without a parseable
yes/no decision are counted as other/invalid. The evidence-probe scorer first
tries to parse JSON and then falls back to labeled fields. For wrong-evidence
rows, both \texttt{wrong\_evidence} and \texttt{contradicted} count as
rejecting the candidate evidence, because both labels indicate that the
evidence should not be accepted as support for the claim. This scoring rule is
why the paper reports parser quality next to the main probe metrics: an
evidence-verification gain should not come from malformed output or a fragile
parser recovery.

\fi

\end{document}
